{\scriptsize
\begin{flushleft}
\end{flushleft}

        \begin{tabularx}{0.49\linewidth}{|p{0.15\columnwidth}|X|}
            \hline 
            \textbf{Notions et contenus} & \textbf{Capacités exigibles} \\
            \hline 
            \multicolumn{2}{|l|}{\textbf{5.4. Équations de Maxwell}} \\
            \hline 
            \multicolumn{2}{|l|}{\textbf{5.4.1. Postulats de l'électromagnétisme }} \\
            \hline
            Force de Lorentz. Équations locales de Maxwell. Formes intégrales. &  Utiliser les équations de Maxwell sous forme locale ou intégrale.
            Relier l'équation de Maxwell-Faraday et la loi de Faraday.
            Établir l'équation locale de la conservation de la charge à partir des équations de Maxwell.
            Utiliser une méthode de superposition. \\ 
            \hline
            \multicolumn{2}{|l|}{\textbf{5.4.2. Aspects énergétiques}} \\
            \hline 
            Vecteur de Poynting. Densité volumique d'énergie électromagnétique. Équation locale de Poynting..  & 
            Utiliser les grandeurs énergétiques pour conduire des bilans d'énergie électromagnétique.
Associer le vecteur de Poynting et lintensité lumineuse utilisée dans le domaine de l'optique. \\
        \hline 
        \end{tabularx}
        \begin{tabularx}{0.49\linewidth}{|p{0.15\columnwidth}|X|}
            \hline 
            \multicolumn{2}{|l|}{\textbf{5.4.3. Approximation des régimes quasi-stationnaires magnétique }} \\
            \hline
            Équations de propagation des champs électrique et magnétique dans le vide. &  Établir les équations de propagation des champs électrique et magnétique dans le vide.
            Expliquer le caractère non instantané des interactions électromagnétiques. \\ 
            \hline
            ARQS magnétique.  & 
            Discuter l'approximation des régimes quasistationnaires.
Simplifier et utiliser les équations de Maxwell et l'équation de conservation de la charge dans l'approximation du régime quasi-stationnaire.
Etendre le domaine de validité des expressions des champs magnétiques obtenues en régime stationnaire. \\
        \hline 
        \end{tabularx}
\vspace{0.5cm}
}
